\begin{notechapter}
\chapter{Distributions and Densities}\label{ch:distributions-densities}
\chaptersubtitle{What a formula written with $p(\cdot)$ means}

\begin{chapterintro}
The book writes $p(x)$, $p(x, y)$ and $p(y\mid x)$ throughout, and most of the time these are densities.
Sometimes no density exists: the empirical distribution of a sample (\cref{cross-entropy:def:empirical}),
the output of a ReLU, a variable that is a deterministic function of another, a point constrained to a
curve. The \emph{distribution} always exists; the \emph{density} is one way of writing it, available only
in some cases. This appendix separates the two notions, and explains in what sense a formula written
with densities remains correct when there is no density: it is then a shorthand for a statement about
probabilities or expectations. You will learn
\begin{itemize}[nosep]
  \item what a distribution is, what a density is, and when the second exists;
  \item densities with respect to other reference measures, which put sums and integrals, and discrete
        and continuous variables, under one notation;
  \item how a density changes with the coordinates, and why a mode is not invariant;
  \item how to read a formula written with densities when there is none: sampling recipes, conditional
        distributions, Bayes' rule with a prior that has no density, density ratios and KL divergences.
\end{itemize}
\prerequisites{integrals and expectations; \cref{ch:learning-problem} (the distribution of $(X, Y)$);
\cref{ch:bayesian-reasoning} (Bayes' theorem) for the last sections.}
\end{chapterintro}

\section{Distribution and density}\label{distributions-densities:sec:basics}

\begin{definition}[Distribution]\label{distributions-densities:def:distribution}
The \emph{distribution} of a random variable $X$ with values in $\R^d$ is the function $P_X$ that gives,
for every (measurable) region $A\subset\R^d$, the probability $P_X(A) = \Prob(X\in A)$.
\end{definition}

The distribution is equivalently described by the expectations $\E[g(X)]$ for all bounded continuous
functions $g$: the indicator $g = \bbone_A$ gives back $P_X(A)$, and indicators can be approximated by
continuous functions. Every random variable has a distribution. This is the object that the book
calls $P$ in \cref{ch:learning-problem}, and the object a sampler produces: drawing $X$ many times and
counting how often it falls in $A$ estimates $P_X(A)$.

\begin{definition}[Density]\label{distributions-densities:def:density}
A \emph{density} of $X$ is a function $p\geq0$ on $\R^d$ such that
\[
  \Prob(X\in A) = \int_Ap(x)\,dx\qquad\text{for every region }A.
\]
\end{definition}

A density is a way of writing the distribution, and a density determines the distribution. Two
densities of the same distribution can differ only on a set of zero volume: changing $p$ at a single
point changes no integral. ``The value of the density at $x$'' is therefore meaningful only for a
continuous version of $p$, and then it is a probability \emph{per unit volume},
\[
  p(x) = \lim_{r\to0}\frac{\Prob(X\in B(x, r))}{\operatorname{vol}(B(x, r))},
\]
where $B(x, r)$ is the ball of radius $r$ around $x$. It is not a probability. It has units (probability
per unit of $x$ in each coordinate), it can exceed $1$ (the uniform density on $[0, \frac12]$ equals $2$),
and the probability of a single point is $0$ whether or not there is a density.

\paragraph{When there is no density.} If $X$ has a density, every region of zero volume has
probability zero, since the integral of $p$ over it is zero. So a distribution that puts positive
probability on a set of zero volume has no density. The converse also holds: if every set of zero
volume has probability zero, a density exists. This is the Radon--Nikodym theorem, whose proof is
beyond the scope of this book. Four cases occur in the book.
\begin{itemize}
  \item \textbf{Discrete values.} The result of a die, a class label $Y$, a word: all the probability is on
        finitely or countably many points.
  \item \textbf{The empirical distribution.} The distribution $\hat p_n$ of a sample puts mass $1/n$ on each
        observed point (\cref{cross-entropy:def:empirical}). It is discrete even when the data are
        continuous.
  \item \textbf{Mixed distributions.} If $Z\sim\mathcal N(0, 1)$, then $Y = \max(0, Z)$, a ReLU unit fed a
        Gaussian pre-activation, has $\Prob(Y = 0) = \frac12$: an atom at $0$. On $(0, \infty)$ its
        probability is spread with the Gaussian density. $Y$ is neither discrete nor has a density
        (\cref{distributions-densities:fig:examples}, left).
  \item \textbf{Fewer degrees of freedom than coordinates.} A point $X = (\cos T, \sin T)$ with $T$ uniform on
        $[0, 2\pi)$ lies on the unit circle, a curve of zero area: $X$ has no density in the plane. The
        same happens to a pair $(X, g(X))$, which lies on the graph of $g$, and to two views of the same
        image when one is a deterministic function of the other.
\end{itemize}

\begin{figure}[ht]
\centering
\begin{tikzpicture}
\begin{axis}[name=relu, width=0.48\linewidth, height=5cm, xmin=-0.6, xmax=3.2, ymin=0, ymax=0.6,
  xlabel={$y$}, axis lines=left, ytick={0, 0.2, 0.4}, tick label style={font=\scriptsize},
  label style={font=\small}, clip=false]
\addplot[thick, bookaccent, domain=0.001:3.2, samples=60] {exp(-x^2/2)/sqrt(2*pi)};
\draw[very thick, bookpitfall, -{Stealth[length=2mm]}] (axis cs:0, 0) -- (axis cs:0, 0.5);
\node[font=\scriptsize, anchor=west, bookpitfall] at (axis cs:0.06, 0.52) {atom $\Prob(Y = 0) = \frac12$};
\node[font=\scriptsize, anchor=west, bookaccent] at (axis cs:0.9, 0.3) {density $\varphi(y)$ on $(0, \infty)$};
\end{axis}
\begin{axis}[at={(relu.east)}, anchor=west, xshift=1.2cm, width=0.44\linewidth, height=5cm,
  axis equal image, xmin=-1.5, xmax=1.5, ymin=-1.5, ymax=1.5, axis lines=none]
\addplot[thick, bookpitfall, domain=0:360, samples=90] ({cos(x)}, {sin(x)});
\addplot[only marks, mark=*, mark size=0.9pt, bookaccent] coordinates {
(0.93,0.27) (0.09,0.85) (0.43,-1.25) (-1.01,-0.39) (1.00,0.23) (-0.99,0.46) (-1.08,0.37) (0.36,0.90)
(-0.25,-0.78) (0.75,0.60) (-1.03,0.88) (-0.88,0.01) (-0.74,0.47) (-0.95,-0.64) (-0.20,-0.79) (0.74,-0.36)
(-0.26,1.01) (-0.51,-0.99) (-0.59,-0.94) (-0.08,1.08) (1.03,-0.04) (1.03,-0.20) (-0.18,0.83) (-0.37,0.90)
(0.86,-0.60) (-0.48,-0.29) (-0.76,-0.13) (0.09,-1.08) (1.06,-0.15) (-0.09,-0.80) (-0.90,0.56) (0.72,0.55)
(-0.44,-0.54) (0.88,-0.30) (0.29,1.23) (-0.57,-0.52) (-0.46,0.93) (-0.17,-0.77) (-0.08,-1.03) (0.13,1.05)
(0.38,-1.02) (-0.48,-0.47) (-0.45,-1.00) (0.25,-1.10) (-0.82,0.56) (0.06,-0.95) (0.74,-0.67) (0.57,0.53)
(0.60,-0.93) (-0.86,0.50) (-1.04,0.26) (0.55,0.77) (-0.20,-0.85) (-0.01,0.65) (0.82,-0.80) (-0.14,1.11)
(-0.76,-0.22) (-0.78,0.83) (-0.80,-0.67) (0.45,0.86)};
\end{axis}
\end{tikzpicture}
\caption{Left: the distribution of $Y = \max(0, Z)$ with $Z\sim\mathcal N(0, 1)$, an atom of mass $\frac12$
at $0$ (arrow) and the Gaussian density on $(0, \infty)$. It has no density on $\R$. Right: $X$ uniform on the
unit circle (red) has no density in the plane; $Z = X + V$ with Gaussian noise $V$ ($\sigma = 0.15$, $60$
draws in blue) has one.}
\label{distributions-densities:fig:examples}
\end{figure}

\section{Densities with respect to other reference measures}\label{distributions-densities:sec:reference}

The volume in \cref{distributions-densities:def:density} is not the only possible yardstick. Replace it
by the \emph{counting measure}, which gives to each set of integers its number of elements: integrals
against it are sums, and a density of a discrete $Y$ with respect to it is the probability mass function,
$\Prob(Y\in A) = \sum_{y\in A}p(y)$. In general, a density of $X$ with respect to a reference measure $\mu$ is a
function $p\geq0$ with $\Prob(X\in A) = \int_Ap\,d\mu$ for every $A$, and then $\E[g(X)] = \int gp\,d\mu$. One
notation covers sums and integrals.

\paragraph{Mixed pairs.} In classification, $X\in\R^D$ is continuous and $Y\in\{1, \dots, K\}$ is a label.
The pair $(X, Y)$ has no density on $\R^{D+1}$ (all its probability is on $K$ parallel copies of $\R^D$), but
it has one with respect to ``volume in $x$ times counting in $y$'':
\[
  \Prob(X\in A, Y = k) = \int_Ap(x, k)\,dx,\qquad p(x, k) = \Prob(Y = k)\,p(x\mid Y = k).
\]
This is the sense of $p(x, y)$ in \cref{ch:learning-problem,ch:generative-discriminative}, and the
reason why sums over $y$ and integrals over $x$ appear side by side in the same formula.

\paragraph{The ReLU output.} With the reference measure ``length on $(0, \infty)$ plus a unit atom at $0$'',
$Y = \max(0, Z)$ has the density $p(0) = \frac12$ and $p(y) = \varphi(y)$ for $y > 0$, where $\varphi$ is the
standard Gaussian density. Then $\E[Y] = 0\cdot\frac12 + \int_0^\infty y\varphi(y)\,dy = 1/\sqrt{2\pi}\approx0.399$.

\begin{pitfall}[Density values depend on the reference measure]
The value $p(0) = \frac12$ above is a probability; the value $p(y) = \varphi(y)$ for $y > 0$ is a probability
per unit length. Comparing them is meaningless, and so is adding the log-likelihood of a discrete
observation to that of a continuous one as if they were on the same scale: the second changes when the
continuous variable is measured in other units, the first does not. For the same reason a likelihood
can be unbounded. A Gaussian mixture with one component centred on a data point and a variance
tending to $0$ has a density at that point tending to infinity, although the probability of any region
stays at most $1$.
\end{pitfall}

\section{Change of variables}\label{distributions-densities:sec:change}

A smooth bijection $T$ of $\R^d$ carries the distribution of $X$ to that of $Y = T(X)$:
$\Prob(Y\in A) = \Prob(X\in T^{-1}(A))$, with no condition. The density, when there is one, changes by the
Jacobian factor that converts volumes:
\[
  p_Y(y) = p_X\bigl(T^{-1}(y)\bigr)\,\abs{\det J_{T^{-1}}(y)} .
\]
The distribution is transported unchanged; its density is not, because ``probability per unit volume''
depends on how volume is measured in the new coordinates.

\begin{example}[The mode moves]\label{distributions-densities:ex:lognormal}
Let $X\sim\mathcal N(0, 1)$ and $Y = e^X$. Then $p_Y(y) = \varphi(\log y)/y$ for $y > 0$, and setting the
derivative of $\log p_Y(y) = -\frac12(\log y)^2 - \log y + \text{const}$ to zero gives $\log y = -1$: the mode of
$Y$ is $e^{-1}\approx0.368$, not $e^0 = 1$, the image of the mode of $X$. The median, which is defined by
probabilities, is transported: it is $e^0 = 1$. The mean is $e^{1/2}\approx1.649$.
\end{example}

This is why a MAP estimate depends on the parametrization while posterior probabilities do not
(\cref{bayesian-reasoning:sec:map}). If $T$ maps $\R^d$ into a space of higher dimension, for instance
$x\mapsto(x, g(x))$, the image has no density at all (\cref{distributions-densities:sec:basics}). In the other
direction, adding an independent variable that has a density restores one: $X + V$ has a density as soon as
$V$ has one, whatever $X$ is (\cref{distributions-densities:ex:circle}).

\section{Formulas with densities when there is none}\label{distributions-densities:sec:notation}

A formula written with densities is a statement about probabilities or expectations. When no density
exists, the book keeps the notation and the statement keeps its meaning; what changes is how it is
read. Four readings cover the book.

\paragraph{1. A chain of conditionals is a sampling recipe.} ``$X\sim p(x)$, then $Y\sim p(y\mid x)$'' says
how $(X, Y)$ is produced: draw $X$, then draw $Y$ from a distribution that depends on the value drawn. A
\emph{conditional distribution} of $Y$ given $X$ is exactly this: for each value $x$, a distribution
$P_{Y\mid X = x}$, such that
\[
  \Prob(X\in A, Y\in B) = \int_AP_{Y\mid X = x}(B)\,dP_X(x)\qquad\text{for all }A, B .
\]
This makes sense for any distributions, with or without densities, and it is the reading of ancestral
sampling (\cref{monte-carlo:sec:hierarchical}) and of the generative models of Part~II onwards. A
deterministic step is the extreme case: ``$Y = g(X)$'' is the conditional distribution with all its mass at
$g(x)$, often written $p(y\mid x) = \delta(y - g(x))$ with the Dirac notation. When all the factors have densities,
the joint density is their product and $p(y\mid x) = p(x, y)/p(x)$; when some do not, the product is the
recipe and the division is not available.

\paragraph{2. Bayes' rule needs a density only for the likelihood.} Let $X$ have any distribution $P_X$, and
let $Z$ given $X = x$ have a density $p(z\mid x)$ in $z$ (a sensor with Gaussian noise, say). Then $Z$ has a
density even when $X$ does not.

\begin{proposition}[Bayes' rule with a general prior]\label{distributions-densities:prop:bayes}
Under these assumptions $Z$ has the density $p(z) = \E[p(z\mid X)] = \int p(z\mid x)\,dP_X(x)$, and for every $z$ with
$p(z) > 0$ a conditional distribution of $X$ given $Z = z$ is
\[
  \Prob(X\in A\mid Z = z) = \frac{\E\bigl[\bbone_A(X)\,p(z\mid X)\bigr]}{p(z)}
  = \frac{1}{p(z)}\int_Ap(z\mid x)\,dP_X(x) .
\]
\end{proposition}

\begin{proof}
For regions $A$ and $B$, the recipe of reading~1 and Fubini's theorem give
\[
  \Prob(X\in A, Z\in B) = \E\Bigl[\bbone_A(X)\int_Bp(z\mid X)\,dz\Bigr]
  = \int_B\E\bigl[\bbone_A(X)\,p(z\mid X)\bigr]\,dz .
\]
With $A = \R^d$ this says that $Z$ has the density $p(z) = \E[p(z\mid X)]$. For general $A$, the integrand is
$p(z)$ times the proposed $\Prob(X\in A\mid Z = z)$ (where $p(z) = 0$ the integrand is $0$ as well), which is what
it means for this formula to be a conditional distribution of $X$ given $Z$ (reading~1 with the roles of
the two variables exchanged).
\end{proof}

When $X$ has a density $p(x)$, the proposition is the familiar
$p(x\mid z) = p(z\mid x)p(x)/p(z)$. In general it says that the posterior has a density \emph{with respect to
the prior}, the ratio $p(z\mid x)/p(z)$: the likelihood reweights the prior, whatever the prior is.

\begin{example}[A point on a circle, observed with noise]\label{distributions-densities:ex:circle}
Let $X$ be uniform on the unit circle, $X = (\cos T, \sin T)$, and $Z = X + V$ with
$V\sim\mathcal N(0, \sigma^2I)$ independent of $X$ (\cref{distributions-densities:fig:examples}, right). $X$ has no
density in the plane, but $p(z\mid x)$ is the Gaussian density, so $Z$ has the density
\[
  p(z) = \frac1{2\pi}\int_0^{2\pi}\frac{1}{2\pi\sigma^2}\,e^{-\norm{z - (\cos t, \sin t)}^2/(2\sigma^2)}\,dt
  = \frac{1}{2\pi\sigma^2}\,e^{-(\norm z^2 + 1)/(2\sigma^2)}\,I_0\bigl(\norm z/\sigma^2\bigr),
\]
where $I_0$ is the modified Bessel function of the first kind. By \cref{distributions-densities:prop:bayes} the
posterior of $X$ reweights the uniform distribution on the circle by
$p(z\mid x)\propto e^{z^\top x/\sigma^2}$: in terms of the angle, it has the density proportional to
$e^{\kappa\cos(t - \alpha)}$, with $\alpha$ the direction of $z$ and $\kappa = \norm z/\sigma^2$, a von Mises
distribution. For $z = (0.8, 0.6)$ and $\sigma = 0.5$: $\kappa = 4$, $p(z)\approx0.132$, and the angle lies within
$0.5$ rad of $\alpha\approx0.644$ with probability $0.66$. The posterior is a distribution on the circle, a
curve: it has no density in the plane either, only a density along the circle.
\end{example}

\paragraph{3. A density ratio needs the joint to be spread like the product.} The ratio
$r(x, y) = p(x, y)/(p(x)p(y))$ of \cref{ch:contrastive-learning}, and the mutual information
$I(X; Y) = \E[\log r(X, Y)]$, require the joint distribution to put zero probability wherever the
product of the marginals does. Two continuous views with $Y = g(X)$ violate this: the joint lies on the
graph of $g$, which has zero probability under the product when $Y$ has no atoms. Then $r$ does not exist
and $I(X; Y) = \infty$. The InfoNCE bound $\log N - \mathcal L_N\leq I(X;Y)$ is still true, but says nothing.

\paragraph{4. A KL divergence needs the same condition.} In general
$\mathrm{KL}(P\Vert Q) = \E_P\bigl[\log\frac{dP}{dQ}\bigr]$, where $\frac{dP}{dQ}$ is the density of $P$ with respect to
$Q$. It exists when every set of $Q$-probability zero has $P$-probability zero, and otherwise
$\mathrm{KL}(P\Vert Q) = \infty$. With densities $p$ and $q$, $\frac{dP}{dQ} = p/q$ and the usual formula follows. For
continuous data and a model $q$ with a density, both $\mathrm{KL}(\hat p_n\Vert q)$ and $\mathrm{KL}(q\Vert\hat p_n)$
are infinite: $\hat p_n$ puts mass on points, which have $q$-probability zero, and $q$ puts mass where
$\hat p_n$ has none (\cref{cross-entropy:sec:forward-reverse}). This is why data are compared with a model
through the cross-entropy $\E_{\hat p_n}[-\log q]$, the average negative log-likelihood, which is finite; with
discrete labels, the conditional KL divergences of \cref{cross-entropy:thm:mle} are finite too.

\begin{keyidea}
A distribution always exists; a density exists only when no set of zero volume carries probability, and
its values depend on the coordinates and on the reference measure. When the book writes $p(\cdot)$ for a
distribution without a density, the formula is a sampling recipe or a statement about expectations,
and it remains correct in that sense.
\end{keyidea}

\begin{chaptersummary}
\begin{itemize}
  \item The distribution of $X$ gives $\Prob(X\in A)$ for every region $A$; a density is a function whose
        integrals give these probabilities. Its values are probabilities per unit volume, not probabilities.
  \item A density with respect to volume exists exactly when sets of zero volume have probability zero
        (Radon--Nikodym). Discrete variables, empirical distributions, ReLU outputs and points on curves
        have none.
  \item Densities with respect to other reference measures (counting, or mixtures of volume and counting)
        cover discrete and mixed variables, and give a meaning to $p(x, y)$ in classification.
  \item Under a change of variables the distribution is transported, the density is multiplied by a
        Jacobian, and the mode moves; probabilities, medians and expectations do not depend on the coordinates.
  \item A chain of conditional distributions is a sampling recipe, valid with or without densities. Bayes'
        rule needs a density only for the likelihood: the posterior reweights the prior by $p(z\mid x)/p(z)$.
  \item Density ratios, mutual information and KL divergences require one distribution to be spread like
        the other; otherwise they are infinite.
\end{itemize}
\end{chaptersummary}

\paragraph{Further reading.} Williams, \emph{Probability with Martingales}, 1991, introduces measure-theoretic
probability briefly and clearly; Capiński and Kopp, \emph{Measure, Integral and Probability}, 2nd ed., 2004,
cover the Radon--Nikodym theorem and conditional expectations at an introductory level; Pollard,
\emph{A User's Guide to Measure Theoretic Probability}, 2002, treats conditional distributions and densities
with respect to general measures in detail.
\end{notechapter}
